\documentclass[a4paper]{article}

\usepackage[spanish]{babel}
\usepackage{graphicx}
\usepackage{amsmath, amssymb}
\usepackage[margin=2cm]{geometry}
\usepackage{fancyhdr}
\usepackage{enumerate}
\usepackage[shortlabels]{enumitem}
\usepackage{parskip}
\usepackage[most]{tcolorbox}
\usepackage[hidelinks]{hyperref}
\usepackage{float}
\usepackage{booktabs}
\usepackage{mathrsfs}
\usepackage{tikz}

% cabecera
\pagestyle{fancy}
\fancyhead[l]{Astroestadística}
\fancyhead[c]{Quiz 1 ???}
\fancyhead[r]{\today}
\fancyfoot[c]{\thepage}
\renewcommand{\headrulewidth}{0.2pt} % linea horizontal


\begin{document}

\section{Quiz de Probabilidad y Estadística}

Seleccione la opción correcta en cada pregunta.

\begin{enumerate}

\item Si dos eventos A y B son independientes, entonces:
\begin{enumerate}[a)]
    \item $P(A \cap B) = P(A) + P(B)$
    \item $P(A|B) = P(A)$
    \item $P(A|B) = P(B|A)$
    \item $P(A \cup B) = 1$
\end{enumerate}

\item ¿Cuál es la diferencia principal entre eventos mutuamente excluyentes e independientes?
\begin{enumerate}[a)]
    \item En ambos casos la intersección es vacía.
    \item Los excluyentes no pueden ocurrir juntos; los independientes no afectan sus probabilidades.
    \item Los excluyentes dependen uno del otro; los independientes son incompatibles.
    \item No hay diferencia práctica entre ambos.
\end{enumerate}

\item En técnicas de conteo, ¿cuándo se usa una permutación en lugar de una combinación?
\begin{enumerate}[a)]
    \item Cuando el orden no importa.
    \item Cuando se reemplazan elementos.
    \item Cuando el orden sí importa.
    \item Cuando hay repeticiones idénticas.
\end{enumerate}

\item ¿Cuál es una condición necesaria para aplicar el teorema de Bayes?
\begin{enumerate}[a)]
    \item Que los eventos sean independientes.
    \item Que el espacio muestral sea infinito.
    \item Que las probabilidades sean mayores que cero.
    \item Que haya al menos dos eventos mutuamente excluyentes que cubran todo el espacio.
\end{enumerate}

\item Una empresa tiene tres fábricas: A, B y C. El 40\% de los productos vienen de la fábrica A, el 35\% de la B y el 25\% de la C. Las tasas de defectos en cada fábrica son 2\% para A, 3\% para B y 5\% para C. ¿Qué representa la expresión siguiente?
$$P(D) = 0.02 \times 0.40 + 0.03 \times 0.35 + 0.05 \times 0.25$$
\begin{enumerate}[a)]
    \item La probabilidad de que un producto provenga de la fábrica A.
    \item La probabilidad de que un producto sea defectuoso, sabiendo que viene de la fábrica A.
    \item La probabilidad de que un producto cualquiera sea defectuoso.
    \item La probabilidad de que un producto defectuoso haya sido fabricado en la fábrica B.
\end{enumerate}

\end{enumerate}

\bibliographystyle{plainurl}
\bibliography{bibliografia} 
\end{document}
